\SourceLeaf{26}
10. Wenn $\xi$ eine Primfunction mod. $(1+x+xx)$ und ihr Determinant $=D$,
so ist allgemein

\[
\zeta^D \equiv \zeta \quad(\mathrm{mod.}\ 1+x+xx,\ \xi).
\]

11. Quadratische, cubische Reste Modulus $\ldots$

12. Ist $P$ cubischer Rest von $\xi$, so ist:

\[
P^{\frac{1}{3}(D-1)} \equiv 1 \quad(\mathrm{mod.}\ 1+x+xx,\ \xi).
\]



\Needspace{8\baselineskip}\begingroup\small
\begin{center}
\textbf{BEMERKUNGEN UBER DIE ENTWICKELUNGEN »ZUR THEORIE DER CUBISCHEN} 
\textbf{RESTE«.}
\end{center}
\noindent\hspace*{\parindent}\textsc{Gauss} hat hier den Versuch gemacht, seinen dritten Beweis des
quadratischen Reciprocitätsgesetzes auf die cubischen Reste zu
übertragen. Der erste Ansatz bezieht sich allein auf das Gebiet der
rationalen ganzen Zahlen und benutzt folgenden Grundgedanken:

Ist die Primzahl $p \equiv 1 \quad(\mathrm{mod.}\ 3)$, so hat die Congruenz $x^3 \equiv 1
\quad(\mathrm{mod.}\ p)$ drei verschiedene Wurzeln, unter denen eine von $1$ verschiedene
$f$ heisse. In dem System aller $p-1$ modulo $p$ incongruenten und gegen
$p$ primen ganzen Zahlen bilden dann die drei Zahlen $1, f, f^2$
gegenüber Multiplication eine Gruppe des Index $\dfrac{p-1}{3}$, für
welche wir auf irgend eine der mannigfachen möglichen Arten ein
Repräsentantensystem $S$, bestehend aus den Zahlen $\alpha_1, \alpha_2,
\ldots, \alpha_{\frac{p-1}{3}}$, ausgewählt denken. Bei diesen Zahlen ist
charakteristisch, dass die Congruenz $\alpha_i \equiv f^{\pm1}\alpha_k
\quad(\mathrm{mod.}\ p)$ für keine zwei Indices $i, k$ bestehen kann. Definirt man
übrigens die Systeme $S'$ und $S''$ zu je $\dfrac{p-1}{3}$ Zahlen
$\beta, \gamma$ durch

\[
\beta_i \equiv f\alpha_i, \qquad \gamma_i \equiv f^2\alpha_i \quad(\mathrm{mod.}\ p),
\]
\noindent so erschöpfen die $(p-1)$ Zahlen $\alpha, \beta, \gamma$ gerade das
gesammte Restsystem\hfill\\
$1, 2, \ldots, {p-1} \quad(\mathrm{mod.}\ p)$.

Ist $a$ eine beliebige durch $p$ nicht theilbare Zahl, so ist

\[
a^{\frac{p-1}{3}} \equiv f^{\lambda} \equiv \left[\frac{a}{p}\right] \quad(\mathrm{mod.}\ p)
\]
\noindent der »cubische Charakter« von $a$ in Bezug auf $p$, und es ist hierbei
$\lambda = 0, 1$ oder $2$.

Um einen Satz über diesen cubischen Charakter zu gewinnen, bilde man die
$\dfrac{p-1}{3}$ Producte $a\alpha_1, a\alpha_2, \ldots
a\alpha_{\frac{p-1}{3}}$. Diese Zahlen bilden offenbar wieder ein
Repräsentantensystem, da aus der Congruenz $a\alpha_i \equiv
f^{\pm1}a\alpha_k$ sofort $\alpha_i \equiv f^{\pm1}\alpha_k \quad(\mathrm{mod.}\ p)$
folgen würde. Gehören nun $\mu$ unter den Zahlen $a\alpha_i$ dem System
$S'$, $\nu$ dem System $S''$ an, so ist offenbar ihr Product mod. $p$ mit
$f^{\mu+2\nu}. \alpha_1\alpha_2\ldots\alpha_{\frac{p-1}{3}}$
congruent; man gewinnt also als Ergebniss:

\[
\left[\frac{a}{p}\right] \equiv a^{\frac{p-1}{3}} \equiv f^{\mu+2\nu} \quad(\mathrm{mod.}\ p),
\]
\noindent welches das genaue Analogon zu dem bekannten Lemma ist, das \textsc{Gauss} seinem
dritten Beweise des quadratischen Reciprocitätsgesetzes zu Grunde legte.

